\documentclass[12pt,letterpaper]{article}

\usepackage[utf8]{inputenc}
\usepackage[T1]{fontenc}
\usepackage[spanish]{babel}
\usepackage{geometry}
\usepackage{graphicx}
\usepackage{float}
\usepackage{booktabs}
\usepackage{longtable}
\usepackage{array}
\usepackage{amsmath}
\usepackage{hyperref}
\usepackage{xcolor}

\geometry{margin=2.5cm}
\setlength{\parskip}{0.55em}
\setlength{\parindent}{0pt}
\renewcommand{\arraystretch}{1.15}

\hypersetup{
  colorlinks=true,
  linkcolor=black,
  urlcolor=blue
}

\newcommand{\grafico}[2]{
  \begin{figure}[H]
    \centering
    \includegraphics[width=0.92\textwidth]{\detokenize{#1}}
    \caption{#2}
  \end{figure}
}

\begin{document}

\begin{titlepage}
  \centering
  {\Large \textbf{Taller 1. Informe de Análisis de Componentes Principales}\par}
  \vspace{1.5cm}
  {\large Análisis Multivariado\par}
  \vspace{0.4cm}
  {\large Maestría en Inteligencia Artificial y Ciencia de Datos\par}
  \vspace{1.2cm}
  {\large \textbf{Tema:} Análisis de Componentes Principales aplicado a calidad del agua del Río Cauca\par}
  \vspace{0.4cm}
  {\large \textbf{Fuente de datos:} Base de datos/Calidad\_del\_agua\_del\_Rio\_Cauca\_20260419.csv\par}
  \vfill
  {\large Fecha de entrega: Mayo 2 de 2026\par}
\end{titlepage}

\tableofcontents
\newpage

\section*{Resumen ejecutivo}
\addcontentsline{toc}{section}{Resumen ejecutivo}

El Análisis de Componentes Principales aplicado a siete variables físico-químicas del Río Cauca permitió resumir la información en tres componentes con valor propio mayor que 1. Estas tres dimensiones explican el \textbf{77,15 \%} de la variabilidad total. La primera dimensión representa un gradiente de carga contaminante y particulada, dominado por sólidos suspendidos, turbiedad, DBO, DQO y fósforo total. La segunda dimensión contrasta oxígeno disuelto y conductividad eléctrica, por lo que aporta una lectura ambiental complementaria.

Los resultados muestran que sí es posible construir un índice, pero debe interpretarse como \textbf{índice parcial de carga contaminante/particulada}, no como índice global de calidad del agua. También se identificaron observaciones extremas en años y estaciones específicas, entre ellas JUANCHITO 1993 y varias estaciones en 2010, que deben revisarse como posibles eventos ambientales particulares y no eliminarse automáticamente.

\section{Introducción}

El presente informe desarrolla una aplicación técnica de Análisis de Componentes Principales (ACP) sobre una base de datos real de calidad del agua del Río Cauca. El objetivo es resumir la estructura multivariada de un conjunto de variables físico-químicas, identificar relaciones entre indicadores ambientales, observar patrones entre estaciones de monitoreo y evaluar si es posible construir un índice sintético a partir de los componentes principales.

El ACP es pertinente para este caso porque las variables de calidad del agua se encuentran medidas en distintas unidades y pueden presentar correlaciones entre sí. La técnica permite transformar el conjunto original de variables en nuevas dimensiones no correlacionadas, denominadas componentes principales, que concentran la mayor parte posible de la variabilidad total.

\section{Datos y preparación metodológica}

La base original contiene \textbf{2368 registros y 56 variables}. Como primer paso, se realizó una limpieza de nombres de columnas, conversión de variables numéricas y normalización de formatos atípicos, incluyendo separador decimal con coma y notación científica escrita como \texttt{*10E}.

Se construyó el identificador \texttt{anio\_estacion}, combinando el año de muestreo con la estación. Esta decisión permite interpretar cada individuo como una observación espacio-temporal, evitando mezclar mediciones de diferentes años en una misma estación.

Durante la auditoría de valores faltantes se identificaron variables con ausencia total de información, por ejemplo hierro disuelto, manganeso disuelto, sodio disuelto, potasio disuelto, cobre disuelto, zinc disuelto, cianuros, fluoruros, sulfuros y sílice. Por esta razón, el ACP se concentró en variables con disponibilidad y relevancia ambiental.

\begin{table}[H]
\centering
\caption{Variables activas seleccionadas}
\begin{tabular}{p{0.36\textwidth} p{0.54\textwidth}}
\toprule
\textbf{Variable} & \textbf{Interpretación ambiental} \\
\midrule
Turbiedad & Material particulado y transparencia del agua \\
Sólidos suspendidos totales & Carga de sedimentos y partículas \\
Demanda bioquímica de oxígeno (DBO) & Materia orgánica biodegradable \\
Demanda química de oxígeno (DQO) & Carga orgánica y química oxidable \\
Oxígeno disuelto & Condición de oxigenación del agua \\
Conductividad eléctrica & Presencia de sales e iones disueltos \\
Fósforo total & Nutrientes asociados a eutrofización \\
\bottomrule
\end{tabular}
\end{table}

Todas las variables seleccionadas presentaron asimetría positiva severa, por lo cual se aplicó transformación \texttt{log1p}. Posteriormente se eliminaron registros incompletos mediante eliminación por caso completo. La matriz final del ACP quedó conformada por \textbf{1475 observaciones}, equivalentes al \textbf{62,29 \%} de la base de análisis.

\section{Estadísticas descriptivas, correlaciones y gráficos}

La exploración inicial muestra alta dispersión en las variables seleccionadas. Los coeficientes de variación fueron elevados, especialmente en DBO, fósforo total y oxígeno disuelto. Esto evidencia heterogeneidad importante en la calidad del agua observada a través del tiempo y entre estaciones. Además, todas las variables presentaron asimetría positiva mayor que 2, por lo cual se justificó la transformación logarítmica antes del ACP.

\begin{table}[H]
\centering
\caption{Estadísticas descriptivas}
\resizebox{\textwidth}{!}{
\begin{tabular}{lrrrrrrr}
\toprule
\textbf{Variable} & \textbf{n} & \textbf{Media} & \textbf{Desv. est.} & \textbf{Min} & \textbf{Max} & \textbf{Asimetría} & \textbf{CV \%} \\
\midrule
Turbiedad & 2185 & 128,00 & 169,00 & 1,00 & 1900,00 & 2,90 & 132,00 \\
Sólidos suspendidos totales & 2297 & 161,00 & 199,00 & 1,20 & 2112,00 & 3,13 & 124,00 \\
DBO & 2203 & 5,94 & 19,10 & 0,10 & 427,00 & 12,20 & 322,00 \\
DQO & 2134 & 30,20 & 37,90 & 1,46 & 706,00 & 6,90 & 125,00 \\
Oxígeno disuelto & 2289 & 4,70 & 9,39 & 0,00 & 216,00 & 14,70 & 200,00 \\
Conductividad eléctrica & 2305 & 123,00 & 106,00 & 0,00 & 4259,00 & 28,10 & 85,50 \\
Fósforo total & 1850 & 0,30 & 0,88 & 0,00 & 14,20 & 9,64 & 293,00 \\
\bottomrule
\end{tabular}}
\end{table}

El diagnóstico de valores extremos mediante el criterio de Tukey mostró presencia de atípicos univariados en todas las variables, especialmente en turbiedad, DQO, sólidos suspendidos y DBO. Por esta razón, dichos valores no se eliminaron de forma automática; se redujo su peso relativo mediante \texttt{log1p} y se evaluaron posteriormente en el plano factorial.

\begin{table}[H]
\centering
\caption{Diagnóstico de valores atípicos univariados}
\begin{tabular}{lrl}
\toprule
\textbf{Variable} & \textbf{Atípicos} & \textbf{Acción aplicada} \\
\midrule
Turbiedad & 216 & Transformación log1p \\
Sólidos suspendidos totales & 159 & Transformación log1p \\
DBO & 152 & Transformación log1p \\
DQO & 201 & Transformación log1p \\
Oxígeno disuelto & 35 & Transformación log1p \\
Conductividad eléctrica & 15 & Transformación log1p \\
Fósforo total & 95 & Transformación log1p \\
\bottomrule
\end{tabular}
\end{table}

\grafico{Graficos de resultados/distribuciones fisico-quimica_grupo_A.png}{Distribuciones físico-químicas, grupo A}
\grafico{Graficos de resultados/distribuciones fisico-quimica_grupo_B.png}{Distribuciones físico-químicas, grupo B}

La matriz de correlación transformada indica relaciones relevantes. La asociación más fuerte se presenta entre \textbf{turbiedad y sólidos suspendidos totales} (r = 0,80), lo cual es coherente con la interpretación ambiental de carga particulada. También se observa relación entre \textbf{DBO y fósforo total} (r = 0,60), que puede asociarse a procesos de contaminación orgánica y aporte de nutrientes. Por otra parte, el oxígeno disuelto se relaciona negativamente con la conductividad eléctrica (r = -0,54) y con DQO (r = -0,27), sugiriendo una dimensión diferente vinculada con condiciones de oxigenación frente a carga iónica o química.

\begin{table}[H]
\centering
\caption{Matriz de correlación sobre variables transformadas}
\resizebox{\textwidth}{!}{
\begin{tabular}{lrrrrrrr}
\toprule
\textbf{Variable} & \textbf{Turbiedad} & \textbf{SST} & \textbf{DBO} & \textbf{DQO} & \textbf{OD} & \textbf{Conductividad} & \textbf{Fósforo} \\
\midrule
Turbiedad & 1,00 & 0,80 & 0,09 & 0,23 & 0,15 & -0,06 & 0,10 \\
SST & 0,80 & 1,00 & 0,21 & 0,38 & 0,10 & 0,07 & 0,23 \\
DBO & 0,09 & 0,21 & 1,00 & 0,15 & -0,20 & 0,23 & 0,60 \\
DQO & 0,23 & 0,38 & 0,15 & 1,00 & -0,27 & 0,24 & -0,01 \\
OD & 0,15 & 0,10 & -0,20 & -0,27 & 1,00 & -0,54 & 0,16 \\
Conductividad & -0,06 & 0,07 & 0,23 & 0,24 & -0,54 & 1,00 & -0,01 \\
Fósforo & 0,10 & 0,23 & 0,60 & -0,01 & 0,16 & -0,01 & 1,00 \\
\bottomrule
\end{tabular}}
\end{table}

\grafico{Graficos de resultados/Matriz de correlacion: Espacion de variable.png}{Matriz de correlación}

\section{Análisis de Componentes Principales}

El ACP se realizó sobre variables transformadas y estandarizadas. La estandarización es necesaria porque las variables están expresadas en unidades diferentes: UNT, mg/L y microS/cm. Las variables \texttt{anio\_estacion} y \texttt{estaciones} se trataron como cualitativas suplementarias para apoyar la interpretación, sin influir en la construcción de los ejes factoriales.

\begin{table}[H]
\centering
\caption{Valores propios e inercia explicada}
\begin{tabular}{lrrr}
\toprule
\textbf{Componente} & \textbf{Valor propio} & \textbf{\% varianza} & \textbf{\% acumulado} \\
\midrule
Dim 1 & 2,210 & 31,568 & 31,568 \\
Dim 2 & 1,780 & 25,431 & 56,999 \\
Dim 3 & 1,410 & 20,148 & 77,147 \\
Dim 4 & 0,670 & 9,565 & 86,712 \\
Dim 5 & 0,466 & 6,659 & 93,371 \\
Dim 6 & 0,296 & 4,224 & 97,594 \\
Dim 7 & 0,168 & 2,406 & 100,000 \\
\bottomrule
\end{tabular}
\end{table}

Con el criterio de Kaiser, se retienen \textbf{tres componentes} con valores propios mayores que 1. Estos tres componentes explican el \textbf{77,15 \%} de la variabilidad total. El plano factorial 1-2 explica \textbf{57,00 \%}, porcentaje suficiente para una primera lectura gráfica, aunque no agota completamente la estructura multivariada.

\grafico{Graficos de resultados/Sedimentacion de varianza.png}{Sedimentación de la varianza}

\section{Nube de individuos}

La nube de individuos permite observar cómo se distribuyen las observaciones espacio-temporales en el plano factorial formado por las dimensiones 1 y 2. Los puntos más alejados del origen representan observaciones con perfiles físico-químicos más diferenciados respecto al comportamiento promedio.

\begin{table}[H]
\centering
\caption{Individuos mejor representados en el plano 1-2}
\begin{tabular}{lrrr}
\toprule
\textbf{Individuo} & \textbf{Dim 1} & \textbf{Dim 2} & \textbf{Cos2 plano} \\
\midrule
2013\_ANTES SUAREZ & -2,145 & 1,375 & 0,992 \\
2011\_PASO DEL COMERCIO & 1,929 & 1,695 & 0,985 \\
1997\_PUENTE HORMIGUERO & -2,190 & 1,341 & 0,978 \\
1998\_ANTES RIO OVEJAS & -4,093 & -0,461 & 0,976 \\
2000\_PASO DE LA BOLSA & -1,714 & 1,469 & 0,972 \\
\bottomrule
\end{tabular}
\end{table}

Estos casos tienen una representación confiable en el plano 1-2, por lo cual su posición gráfica puede interpretarse con mayor seguridad.

\grafico{Graficos de resultados/Nube de individuos.png}{Nube de individuos}

\section{Círculo de correlaciones}

El círculo de correlaciones muestra la relación entre las variables originales y los componentes principales. En la Dim 1 predominan variables asociadas a carga contaminante y particulada: sólidos suspendidos totales, turbiedad, DBO, DQO y fósforo total. En la Dim 2 se destacan principalmente oxígeno disuelto y conductividad eléctrica, pero con signos opuestos.

\begin{table}[H]
\centering
\caption{Interpretación de los ejes factoriales}
\begin{tabular}{lll}
\toprule
\textbf{Dimensión} & \textbf{Variables principales} & \textbf{Interpretación} \\
\midrule
Dim 1 & Sólidos suspendidos, turbiedad, DBO, DQO, fósforo & Gradiente de carga contaminante, orgánica y particulada \\
Dim 2 & Oxígeno disuelto (+), conductividad eléctrica (-) & Contraste entre oxigenación e influencia iónica/conductiva \\
\bottomrule
\end{tabular}
\end{table}

\grafico{Graficos de resultados/Circulos de correlaciones Plano 1-2.png}{Círculo de correlaciones}

\section{Biplot y datos atípicos}

El biplot integra individuos y variables en una misma representación. Las observaciones ubicadas en la dirección de una variable tienden a presentar valores relativamente altos en esa variable. En el lado positivo de la Dim 1 se ubican observaciones asociadas a mayor carga particulada y orgánica; en la Dim 2 se observa el contraste entre oxígeno disuelto y conductividad eléctrica.

Para identificar posibles datos atípicos se utilizó la distancia al origen en el plano factorial. Los diez casos más alejados fueron:

\begin{table}[H]
\centering
\caption{Posibles atípicos factoriales}
\resizebox{\textwidth}{!}{
\begin{tabular}{llrrrr}
\toprule
\textbf{Individuo} & \textbf{Estación} & \textbf{Dim 1} & \textbf{Dim 2} & \textbf{Distancia} & \textbf{Cos2 plano} \\
\midrule
1993\_JUANCHITO & JUANCHITO & -3,216 & 6,571 & 53,519 & 0,393 \\
2010\_MEDIACANOA & MEDIACANOA & 4,819 & 4,587 & 44,265 & 0,329 \\
2010\_PASO DEL COMERCIO & PASO DEL COMERCIO & 4,734 & 4,277 & 40,700 & 0,317 \\
2010\_PUERTO ISAACS & PUERTO ISAACS & 4,691 & 4,267 & 40,219 & 0,313 \\
2010\_JUANCHITO & JUANCHITO & 3,500 & 5,149 & 38,761 & 0,297 \\
2010\_YOTOCO & YOTOCO & 4,478 & 4,081 & 36,700 & 0,274 \\
1990\_ANTES INTERCEPTOR SUR & ANTES INTERCEPTOR SUR & 4,575 & -3,844 & 35,708 & 0,740 \\
1994\_PASO DEL COMERCIO & PASO DEL COMERCIO & 5,086 & -2,800 & 33,702 & 0,753 \\
2010\_ANACARO & ANACARO & 4,887 & 3,125 & 33,653 & 0,240 \\
1992\_ANTES RIO TIMBA & ANTES RIO TIMBA & -4,925 & 3,058 & 33,600 & 0,554 \\
\bottomrule
\end{tabular}}
\end{table}

Estos registros no deben eliminarse automáticamente, ya que pueden corresponder a eventos reales de contaminación, condiciones hidrológicas particulares o cambios temporales relevantes. Deben interpretarse como observaciones extremas desde el punto de vista factorial.

\grafico{Graficos de resultados/Relacion Dual individuos - variables.png}{Biplot: relación dual individuos-variables}

\section{Síntesis con contribuciones y cosenos cuadrados}

El aporte promedio esperado por variable en cada dimensión es 100/7 = \textbf{14,29 \%}. En la Dim 1, las mayores contribuciones corresponden a sólidos suspendidos totales, turbiedad y DBO. En conjunto, estas variables sostienen la interpretación de la Dim 1 como un gradiente de carga contaminante y particulada.

\begin{table}[H]
\centering
\caption{Contribuciones y cosenos cuadrados en la Dimensión 1}
\begin{tabular}{lrrr}
\toprule
\textbf{Variable} & \textbf{Coord. Dim 1} & \textbf{Contrib. Dim 1} & \textbf{Cos2 Dim 1} \\
\midrule
Sólidos suspendidos totales & 0,846 & 32,358 & 0,715 \\
Turbiedad & 0,720 & 23,453 & 0,518 \\
DBO & 0,589 & 15,712 & 0,347 \\
DQO & 0,547 & 13,519 & 0,299 \\
Fósforo total & 0,491 & 10,913 & 0,241 \\
Conductividad eléctrica & 0,277 & 3,477 & 0,077 \\
Oxígeno disuelto & -0,112 & 0,568 & 0,013 \\
\bottomrule
\end{tabular}
\end{table}

En la Dim 2, las variables dominantes son oxígeno disuelto y conductividad eléctrica. La oposición de signos sugiere que este componente expresa una condición ambiental distinta a la Dim 1.

\begin{table}[H]
\centering
\caption{Contribuciones y cosenos cuadrados en la Dimensión 2}
\begin{tabular}{lrrr}
\toprule
\textbf{Variable} & \textbf{Coord. Dim 2} & \textbf{Contrib. Dim 2} & \textbf{Cos2 Dim 2} \\
\midrule
Oxígeno disuelto & 0,854 & 41,015 & 0,730 \\
Conductividad eléctrica & -0,766 & 32,983 & 0,587 \\
Turbiedad & 0,430 & 10,381 & 0,185 \\
DQO & -0,309 & 5,355 & 0,095 \\
Sólidos suspendidos totales & 0,307 & 5,288 & 0,094 \\
DBO & -0,267 & 3,992 & 0,071 \\
Fósforo total & 0,133 & 0,986 & 0,018 \\
\bottomrule
\end{tabular}
\end{table}

En síntesis, los cosenos cuadrados indican que los sólidos suspendidos y la turbiedad están bien representados por la Dim 1, mientras que oxígeno disuelto y conductividad eléctrica están mejor representados por la Dim 2. Esta separación respalda que el fenómeno no es unidimensional: existe un eje de carga particulada/orgánica y otro eje de condición físico-química asociada a oxigenación y conductividad.

\section{Construcción de índice}

La Dim 1 explica \textbf{31,57 \%} de la varianza total y presenta cargas positivas en 6 de las 7 variables activas. Además, las variables con mayor contribución en esta dimensión tienen una interpretación ambiental coherente: sólidos suspendidos, turbiedad, DBO, DQO y fósforo. Esta coherencia permite construir un indicador ordenado por la posición de los individuos sobre la Dim 1.

Por lo anterior, \textbf{sí es posible construir un índice parcial basado en la Dim 1}, pero no se recomienda interpretarlo como un índice global de calidad del agua. La razón es que la Dim 2 contiene información ambiental importante, especialmente sobre oxígeno disuelto y conductividad eléctrica, que no queda completamente representada por la Dim 1. En consecuencia, el índice se interpreta como \textbf{índice parcial de carga contaminante/particulada}: valores altos indican mayor carga relativa, no necesariamente peor calidad integral en todos los criterios ambientales.

El índice construido se normalizó en escala 0-100 a partir de la coordenada factorial de cada individuo en la Dim 1:

\[
\text{Indice\_PC1\_0\_100} = 100 \times \frac{PC1 - \min(PC1)}{\max(PC1) - \min(PC1)}
\]

Valores cercanos a 100 indican mayor carga contaminante/particulada según la estructura capturada por la Dim 1.

\begin{table}[H]
\centering
\caption{Mayores valores individuales del índice parcial}
\resizebox{\textwidth}{!}{
\begin{tabular}{llrr}
\toprule
\textbf{Individuo} & \textbf{Estación} & \textbf{PC1} & \textbf{Índice} \\
\midrule
2010\_LA VICTORIA & LA VICTORIA & 5,104 & 100,000 \\
1994\_PASO DEL COMERCIO & PASO DEL COMERCIO & 5,086 & 99,825 \\
2010\_ANACARO & ANACARO & 4,887 & 97,906 \\
2010\_MEDIACANOA & MEDIACANOA & 4,819 & 97,245 \\
2010\_PASO DEL COMERCIO & PASO DEL COMERCIO & 4,734 & 96,422 \\
2010\_PUERTO ISAACS & PUERTO ISAACS & 4,691 & 96,009 \\
1990\_ANTES INTERCEPTOR SUR & ANTES INTERCEPTOR SUR & 4,575 & 94,881 \\
2010\_YOTOCO & YOTOCO & 4,478 & 93,942 \\
2010\_PUENTE LA VIRGINIA & PUENTE LA VIRGINIA & 4,458 & 93,750 \\
2010\_PASO DE LA TORRE & PASO DE LA TORRE & 4,315 & 92,371 \\
\bottomrule
\end{tabular}}
\end{table}

\begin{table}[H]
\centering
\caption{Estaciones con mayor índice parcial promedio}
\begin{tabular}{lrrr}
\toprule
\textbf{Estación} & \textbf{n mediciones} & \textbf{Índice promedio} & \textbf{Índice máximo} \\
\midrule
ANACARO & 92 & 57,6 & 97,9 \\
PUENTE LA VIRGINIA & 89 & 57,5 & 93,8 \\
LA VICTORIA & 90 & 57,4 & 100,0 \\
VIJES & 88 & 56,6 & 88,2 \\
PUENTE GUAYABAL & 92 & 56,6 & 88,6 \\
RIOFRIO & 90 & 56,3 & 91,7 \\
YOTOCO & 87 & 56,0 & 93,9 \\
MEDIACANOA & 100 & 55,9 & 97,2 \\
PASO DE LA TORRE & 87 & 53,9 & 92,4 \\
PASO DEL COMERCIO & 96 & 53,4 & 99,8 \\
\bottomrule
\end{tabular}
\end{table}

\section{Hallazgos principales}

\begin{enumerate}
  \item La base de datos presenta alta heterogeneidad y asimetría positiva en las variables físico-químicas analizadas, lo que justificó el uso de transformación logarítmica.
  \item La correlación más fuerte se encontró entre turbiedad y sólidos suspendidos totales, evidenciando una dimensión clara de carga particulada.
  \item El ACP retuvo tres componentes según el criterio de Kaiser, explicando el 77,15 \% de la variabilidad total.
  \item El plano Dim 1-Dim 2 explicó 57,00 \% de la variabilidad, suficiente para interpretar patrones generales de individuos y variables.
  \item La Dim 1 representa principalmente carga contaminante/particulada, mientras que la Dim 2 refleja un contraste entre oxígeno disuelto y conductividad eléctrica.
  \item Se identificaron observaciones extremas, especialmente en estaciones y años como JUANCHITO 1993, MEDIACANOA 2010, PASO DEL COMERCIO 2010 y PUERTO ISAACS 2010.
  \item Es posible construir un índice parcial de carga contaminante basado en PC1, pero no un índice global de calidad del agua, porque PC2 contiene información ambiental relevante no resumida por PC1.
  \item Las estaciones con mayores promedios del índice parcial fueron ANACARO, PUENTE LA VIRGINIA, LA VICTORIA, VIJES y PUENTE GUAYABAL.
\end{enumerate}

\section{Conclusiones}

El Análisis de Componentes Principales permitió resumir adecuadamente la estructura multivariada de las variables físico-químicas del Río Cauca. La primera dimensión sintetiza un gradiente de contaminación asociado a sólidos suspendidos, turbiedad, DBO, DQO y fósforo total. La segunda dimensión complementa la interpretación al representar principalmente el contraste entre oxígeno disuelto y conductividad eléctrica.

Los resultados muestran que la calidad del agua no puede describirse de manera completa con un único eje factorial. Sin embargo, la Dim 1 ofrece una base coherente para construir un índice parcial de carga contaminante/particulada, útil para ordenar observaciones y estaciones según su perfil relativo. Este índice debe interpretarse con cautela y acompañarse de la lectura de la Dim 2 para no perder información ambiental importante.

En términos aplicados, el ACP permitió identificar variables dominantes, estaciones con mayor carga relativa y observaciones extremas que podrían ser objeto de revisión ambiental o análisis complementario.

\end{document}
